% LaTeX Curriculum Vitae Template
%
% Copyright (C) 2004-2009 Jason Blevins <jrblevin@sdf.lonestar.org>
% http://jblevins.org/projects/cv-template/
%
% You may use use this document as a template to create your own CV
% and you may redistribute the source code freely. No attribution is
% required in any resulting documents. I do ask that you please leave
% this notice and the above URL in the source code if you choose to
% redistribute this file.
\documentclass[10pt,letterpaper]{article}
\usepackage{hyperref}
\usepackage{geometry}
\usepackage[T1]{fontenc}
% Comment the following line to use the default Computer Modern font
% instead of the Palatino font provided by the mathpazo package.
% Remove the 'osf' bit if you don't like the old style figures.
\usepackage[sc,osf]{mathpazo}
% In practice, I use the following font packages instead of mathpazo.
% beramono provides a nice fixed-width font. xagaramon uses the
% (commercial) Adobe Garamond font.
%\usepackage[scaled=0.75]{beramono}
%\usepackage[osf]{xagaramon}
% Set your name here
\def\name{Regi Kusumaatmadja}
% Replace this with a link to your CV if you like, or set it empty
% (as in \def\footerlink{}) to remove the link in the footer:
\def\footerlink{https://sites.google.com/view/regi-kusumaatmadja/}
% The following metadata will show up in the PDF properties
\hypersetup{
  colorlinks = true,
  urlcolor = black,
  pdfauthor = {\name},
  pdfkeywords = {economics, econometrics, industrial organization, applied econometrics, innovation,
    applied microeconomics},
  pdftitle = {\name: Curriculum Vitae},
  pdfsubject = {Curriculum Vitae},
  pdfpagemode = UseNone
}
\geometry{
  body={6.5in, 8.5in},
  left=1.0in,
  top=1.25in
}
% Customize page headers
\pagestyle{myheadings}
\markboth{TINBERGEN INSTITUTE}{\name}
\thispagestyle{empty}
% Custom section fonts
\usepackage{sectsty}
\sectionfont{\rmfamily\mdseries\Large}
\subsectionfont{\rmfamily\mdseries\itshape\large}
% Other possible font commands include:
% \ttfamily for teletype,
% \sffamily for sans serif,
% \bfseries for bold,
% \scshape for small caps,
% \normalsize, \large, \Large, \LARGE sizes.
% Don't indent paragraphs.
\setlength\parindent{0em}
% Make lists without bullets
\renewenvironment{itemize}{
  \begin{list}{}{
    \setlength{\leftmargin}{1.5em}
  }
}{
  \end{list}
}
\begin{document}
\centerline{{\huge \name}}
\vspace{0.1in}
% {\small IMPORTANT most agencies like central banks block google drive and dropbox. Use for example Github to host your CV}
\vspace{0.25in}
% Place name at left
% {\huge \name}
% Alternatively, print name centered and bold:
%\centerline{\huge \bf \name}
\vspace{0.25in}
\begin{minipage}{0.59\textwidth}
Department of Economics, VU Amsterdam \\
Email: \href{mailto:r.kusumaatmadja@vu.nl}{r.kusumaatmadja@vu.nl} \\
Website: \href{https://sites.google.com/view/regi-kusumaatmadja}{https://sites.google.com/view/regi-kusumaatmadja}
\end{minipage}
% \hfill
% \begin{minipage}{0.45\textwidth}
% BUSINESS ADDRESS \\
% Department of Economics \\
% VU Amsterdam \\
% De Boelelaan 1105 \\
% 1081 HV Amsterdam \\
% The Netherlands \\
% Tel: 
% \end{minipage}
% \vspace{0.15in}
% Citizenship: Indonesian \hfill Gender: M \hfill Date of Birth: 06/21/1994
\section*{Education}
VU Amsterdam and Tinbergen Institute, Ph.D. in Economics, Expected Completion: 2026. \\
Tinbergen Institute, M.Phil. in Economics, 2021. \\
VU Amsterdam, M.Sc. in Economics, 2018. \\
University of Indonesia, B.Sc. in Physics, 2016 (Best Graduate).
\section*{Letter Writers}
\noindent\begin{tabular}{@{}p{0.48\textwidth}@{\hspace{0.04\textwidth}}p{0.48\textwidth}@{}}
\textbf{prof. dr. Jos\'e Luis Moraga-Gonz\'alez} &
\textbf{prof. dr. Sabien Dobbelaere} \\
VU Amsterdam \& T\'el\'ecom Paris &
VU Amsterdam \\
Tinbergen Institute, CEPR, CESifo &
Tinbergen Institute, IZA \\
Email: \href{mailto:j.l.moragagonzalez@vu.nl}{j.l.moragagonzalez@vu.nl} &
Email: \href{mailto:sabien.dobbelaere@vu.nl}{sabien.dobbelaere@vu.nl} \\[0.3in]
\textbf{prof. dr. Mark J. Roberts} &
\textbf{prof. dr. Eric J. Bartelsman} (placement dir.) \\
Pennsylvania State University &
VU Amsterdam \\
NBER, ZEW &
Tinbergen Institute, IZA \\
Email: \href{mailto:mroberts@psu.edu}{mroberts@psu.edu} &
Email: \href{mailto:e.j.bartelsman@vu.nl}{e.j.bartelsman@vu.nl}
\end{tabular}
% \textbf{prof. dr. Eric J. Bartelsman} (placement director) \\
% Department of Economics, VU Amsterdam \\
% Fellows, Tinbergen Institute and IZA \\
% Email: \href{mailto:e.j.bartelsman@vu.nl}{e.j.bartelsman@vu.nl}
\section*{Visiting Position}
2026.05 - 2026.06, CUNEF Universidad Madrid, invited by \textbf{dr. Costanza Cincotta} \\
2023.01 - 2023.06, the Pennsylvania State University, invited by \textbf{prof. dr. Mark J. Roberts}
\section*{Research and Teaching Areas}
Primary Fields: Industrial Organization, Economics of Innovation \\
Secondary fields: Applied Econometrics, Applied Microeconomics 
% Auto-generated by extract_papers.py. Edit papers.json or the paper titleInput.tex files.
\section*{Research in Progress}
\textbf{Internal and External R\&D: An analysis of costs and benefits}\\[0.35em]
This paper analyzes the costs and benefits of internal and external R\&D activities. Using Dutch production and innovation surveys between 2000 and 2020, focusing on the ICT industry, I document an increasing share of firms performing R\&D and the four-way mix of in-house research, contracted-out research, both, and neither. To rationalize these findings, I build and estimate a dynamic discrete-choice model of R\&D with mode-specific investment costs. The cost of contracted-out R\&D is much higher than the cost of in-house R\&D, which explains the small share of firms that contract out alone. Doing both is cheaper than the sum of the two costs: in-house R\&D offsets almost all of the extra cost of going outside. I read that extra cost as a hold-up problem with the partner. The offset is absorptive capacity: a firm that already does in-house R\&D keeps more of the collaboration, so the two activities cost less together than apart. Productivity gains are similar across modes, so the mix of activities is driven by these costs. A uniform subsidy that does not favor a mode raises the share of R\&D-active firms and firm value by more than a subsidy aimed only at in-house R\&D.\par
\vspace*{0.05in}
\textit{Presentation}: \href{http://asesec.org/jornadas_economia_industrial/programme.html}{JEI - Gran Canaria} (2022), \href{https://sites.google.com/view/phd-evs2020/home?authuser=0}{Ph.D.-EVS} (2022), IO/Applied Micro Brownbag - Penn State (2023), VU Amsterdam Internal Seminar (2023), \href{https://earie.org/programme-rome-2023/}{EARIE - Rome} (2023), Tinbergen Institute Ph.D. Seminar (2023), DKIIS - Valencia (2023), \href{https://econ.la.psu.edu/comparative-analysis-of-enterprise-data-caed/}{CAED} - Penn State (2024), \href{https://economicdynamics.org/sedam_2024/}{SED Meeting} - Barcelona (2024), \href{https://us10.forward-to-friend.com/forward/preview?u=bb253d28cb05d3719c9e0c960&id=c939ef0aa0}{CEPR Workshop: Trade, Geography, and Industrial Organization - Erasmus Rotterdam} (2024, poster), \href{https://sites.google.com/view/apioc-2024}{APIOC - Seoul} (2024), CEPR Symposium - Paris (2024, poster), EEA-ESEM - Bordeaux (2025), CEPR IO Virtual Gathering (2025)\par
\vspace*{0.12in}
\textbf{Entry into Innovation Areas}\\[0.35em]
This paper investigates the multi-area entry decisions of publicly listed firms in the United States electronics industry from 1990 to 2019. I apply Latent Dirichlet Allocation to over 630,000 patent titles and abstracts to endogenously define twenty innovation areas and merge these with market-implied patent valuations. The data reveal two empirical patterns. First, innovation occurs across multiple domains simultaneously. Second, the number of active competitors varies across areas. To rationalize these patterns, I estimate a static simultaneous-move game of firm entry with moment inequalities. First, I find that firms enter areas where expected patent valuations are higher. Second, holding those valuations fixed, additional rivals reduce expected profit. Third, firms also benefit from occupying more areas at once. Fourth, a large per-area cost of occupying an area in a given year, together with rival congestion, still limits how many areas a firm enters. Starting from the areas firms already occupy, and holding rival counts fixed, I find that a 25\% reduction in the yearly per-area cost raises to 22--27\% the share of firms that occupy an extra area in a given year. A 50\% reduction in how costly rivals are raises that share to 15--21\%. In a case study of targeted cost reductions, those partial-equilibrium reallocations raise occupancy in semiconductor areas such as ``Semiconductor Memory Architecture'' and ``Semiconductor Lithography and Fabrication Processes.'' The offset is lower occupancy in other fields, including ``Wireless Base-Station and Mobile Communications.''\par
\vspace*{0.05in}
\textit{Presentation}: VU Amsterdam Internal Seminar (2025)\par
\vspace*{0.12in}
\textbf{The Impact of Government Financial Assistance on the Scope and Direction of Firm Innovation}, with C. Cincotta.\\[0.35em]
Governments spend hundreds of billions of dollars each year to support private-sector innovation, but it is unclear whether this money expands existing research, redirects it toward new technologies, or reallocates innovation away from other firms. To address this question, we develop a two-stage model in which firms first choose an R\&D portfolio and then compete in prices \`a la Salop. Financially constrained firms underinvest, and even unconstrained firms under-provide quality because product-market competition leaves a gap between private and social returns. Liquidity assistance can relax the cash constraint, while the quality gap remains. Cheaper cash can expand how much a recipient patents and how many technology areas it covers. If the firm simply scales up the research it already conducts, the composition of that research does not change. Assistance that is targeted toward particular activities can redirect the portfolio. Nearby rivals may expand through knowledge pooling or contract through business stealing. We take these implications to a panel of U.S.\ public firms from 2000 to 2021, combining Good Jobs First awards with patent data. Scope measures how much a firm patents and how many technology areas it covers. Direction measures whether those patents align with frontier or peer technologies. To estimate causal effects, we first reweight untreated firms so that they match recipients on pre-treatment characteristics, and we then estimate a stacked difference-in-differences. Our model implies that distant rivals provide a better control group for the treated firm's own response. Nearby rivals may themselves be affected by business stealing or knowledge pooling, so comparing the two groups traces how those nearby rivals respond. We implement this design using both industry codes and text-based product-market overlap. Distant rivals provide the preferred reading of the treated firm's own response. Financial assistance expands the scale of recipients' patent portfolios. Relative to distant industry controls, intensive scope rises by about 10.5 percent after an award. Relative to firms with little product-market overlap, the same contrast is about 3 percent. Direction shifts toward the industry's most-cited technologies, and the shift is strongest for cash grants. When nearby means high text-based product-market overlap and distant means little or none, the product-market spills are indistinguishable from zero, so those nearby rivals' patenting scale does not move. When nearby means the same two-digit industry, the close estimate on scale is smaller than the distant estimate, which means industry-code neighbors expand as well, the knowledge-pooling case.\par
\vspace*{0.12in}
\textbf{Mergers and Innovation: An Exploration of Scope and Direction of Innovation}, with S. Dobbelaere and J. L. Moraga-Gonzalez.\\[0.35em]
How do mergers change the scale and composition of firm innovation? We study U.S.\ public mergers from 1980 to 2020 using combined acquirer--target patent portfolios. We measure \emph{scope}, the scale of the joint portfolio, and \emph{direction}, its alignment with highly cited patents. A compact framework isolates three mechanisms. Common ownership internalizes business-stealing externalities in product-market overlap areas and reduces the private return to innovation there. The merger may also pool technological assets, raising the productivity of R\&D where the firms possess related knowledge. Some innovation lines require a fixed cost before they can be operated; pooled assets can make previously inactive lines worth activating. Product-market overlap and technological overlap are distinct: firms may compete for the same customers without sharing knowledge, or share related technologies without competing for the same customers. We compare treated merger pairs to matched control pairs before and after the deal. On the matched sample, joint-portfolio scope is larger after mergers than for matched controls. The difference is substantially larger when pre-merger technologies overlap, especially among non-horizontal deals. We find no strong evidence that activity is reallocated between technology classes shared by both firms and classes unique to one. Joint portfolios show higher alignment with highly cited patents relative to matched controls, although the framework leaves the sign of that change theoretically ambiguous. Acquirer-only R\&D measures overstate post-merger input growth relative to joint-entity measurement.\par
\vspace*{0.05in}
\textit{Presentation}: VU Amsterdam Internal Seminar (2024), Tinbergen Institute Workshop in Economics of Innovation (2025)\par

\section*{Teaching Experience}
2021-2025 TA, VU Amsterdam, Introduction to Econometrics (MSc). Eval: 4.3/5.0 \\
2021-2025 TA, VU Amsterdam, Applied Econometrics (MSc). Eval: 4.7/5.0 \\
2022-2025 Supervisor, VU Amsterdam, Research Project and Thesis (BSc + MSc). \\
2021-2024 Instructor, VU Amsterdam, Math, Statistics, Micro, and R (Refresher MSc). \\
2022 TA, VU Amsterdam, Microeconomics I (BSc). \\
2021-2025 Guest Lecturer, VU Amsterdam, Urban Economic Challenges \& Policies (MSc). 
\section*{Relevant Experience}
2023, EARIE 1st Summer School on Entry Model, LUISS - Rome. \\
2022, Econometric Society's Dynamic Structural Econometrics Summer School, MIT. \\
2019, Research on Productivity, Trade, and Growth Summer School, Tinbergen Institute. \\
2020-2021, RA for Eric J. Bartelsman and Sabien Dobbelaere, VU Amsterdam, MICROPROD project. \\
2018-2019, Research Associate, Institute for Economic \& Social Research, University of Indonesia (LPEM FEB UI).
\section*{Scholarships, Honors, and Awards}
2025, Excellence in Teaching Award, VU Amsterdam. \\
2023, Research Visit Grant, VU Amsterdam (monetary grant). \\
2019-2021, Tinbergen Institute - Full Scholarship (tuition \& living cost). \\
2017, VU Amsterdam Fellowship Programme (tuition). \\
2017, Holland Scholarship Programme (living cost). \\
2016, Best Graduate, Department of Physics, University of Indonesia (monetary prize). \\
2016, Top 1 (one) Percent, Faculty of Mathematics and Natural Sciences, University of Indonesia. \\
2014-2016, XL-Axiata Future Leaders Scholarship (leadership programme \& in-kind benefits). \\
2014-2015, Karya Salemba Empat Scholarship (tuition \& living cost). 
\section*{Publications}
\textbf{Evaluating the congestion-reducing effects of road rationing policy: Evidence from Jakarta's odd-even policy}. \textit{Case Studies on Transport Policy} (2025), with M. Halley Yudhistira, Y. Sofiyandi, and M. Firman Hidayat. (previously circulated as Does Traffic Management Matter? Evaluating Congestion Effect of Odd-Even Policy in Jakarta)

\vspace{0.1in}

\textbf{Theoretical exploration of optical response of $Fe_{3}O_{4}$-reduced graphene oxide nano-particle system within dynamical mean-field theory}. \textit{IOP Conf. Ser.: Mater. Sci. Eng.} (2017), with M.A. Majidi, A.D. Fauzi, W.Y. Phan, A. Taufik, R. Saleh, and A. Rusydi.
% \section*{Research in Progress}
% \textbf{Internal and External R\&D: An analysis of costs and benefits (Job Market Paper)} 
% This paper analyzes the costs and benefits of internal and external R\&D activities. Using Dutch production and innovation surveys between 2000 and 2020 focusing on the IT industry, I document an increasing trend of R\&D activities across industries and present evidence suggesting that internal and external R\&D are complementary. To rationalize these findings, I build and estimate a dynamic discrete choice model of R\&D, which explicitly includes specific investment costs of R\&D. I find that the cost of doing external R\&D is $\approx 4$ times higher than the internal R\&D, reflecting the transaction costs of such contract and explaining the observed small share of external R\&D firms in the data. To mimic the Dutch Tax Incentives for Innovation scheme, I simulate the effect of two types of R\&D subsidization programs. I find that if the government has no preference for any particular R\&D activity, the share of R\&D-active firms increases the most and leads to higher change in the firm's total valuation. On the other hand, if the government prefers the firms only to perform internal R\&D, the share of R\&D-active firms is practically unchanged and leads to lower change in the valuation.
% \vspace{0.05in}
% \textit{Presentation}: \href{http://asesec.org/jornadas_economia_industrial/programme.html}{JEI - Gran Canaria} (2022), \href{https://sites.google.com/view/phd-evs2020/home?authuser=0}{Ph.D.-EVS} (2022), IO/Applied Micro Brownbag - Penn State (2023), VU Amsterdam Internal Seminar (2023), \href{https://earie.org/programme-rome-2023/}{EARIE - Rome} (2023), Tinbergen Institute Ph.D. Seminar (2023), DKIIS - Valencia (2023), \href{https://econ.la.psu.edu/comparative-analysis-of-enterprise-data-caed/}{CAED} - Penn State (2024), \href{https://economicdynamics.org/sedam_2024/}{SED Meeting} - Barcelona (2024), \href{https://us10.forward-to-friend.com/forward/preview?u=bb253d28cb05d3719c9e0c960&id=c939ef0aa0}{CEPR Workshop: Trade, Geography, and Industrial Organization - Erasmus Rotterdam} (2024, poster), \href{https://sites.google.com/view/apioc-2024}{APIOC - Seoul} (2024), CEPR Symposium - Paris (2024, poster), EEA-ESEM - Bordeaux (2025)
% \vspace{0.1in}
% \textbf{Merger and Innovation: An Exploration of Firms' Scope and Direction of Innovation}, with S. Dobbelaere and J. L. Moraga-Gonzalez.
% Mergers are often justified as enhancing innovation through synergies, yet their effects on the scope (magnitude) and direction (composition) of firm innovation remain debated. This paper develops a theoretical framework integrating mergers with multi-project innovation decisions, showing that: mergers can increase, decrease, or leave unchanged scope and direction depending on project substitutability. Empirically, using a version of the matching estimator on U.S. public mergers from 1980 to 2020, we find no robust overall change in innovation scope but improved directional alignment to high-impact benchmarks by 0.02 (on a 0-1 cosine similarity scale). Heterogeneity reveals downsides: high technological overlap reduces scope by 21–60\% and diverges direction by 0.07; inactive acquirer–active target mergers shrink scope by 47\% and diverge by 0.21; both-active cases cut scope by 16\% and diverge by 0.06. These findings highlight mergers' potential to contract research diversity in overlapping or innovation-intensive deals, informing antitrust policy in tech-heavy industries.
% \vspace{0.05in}
% \textit{Presentation}: VU Amsterdam Internal Seminar (2024), Tinbergen Institute Workshop in Economics of Innovation (2025)
% \vspace{0.1in}
% \textbf{Entry into Innovation Areas} 
% This paper investigates the notion of innovation direction as the outcome of firms' entry decisions into specific innovation areas. Using Latent Dirichlet Allocation, a natural language processing technique in topic modeling, on 631,000 U.S. electronics patents (1990–2019), I first optimally classify each patent and firm's innovation portfolio into 20 innovation areas. Key stylized facts reveal persistent multi-domain engagement, with patents and firms averaging 2.86 and 4.68 topics respectively, alongside rising entry in software-intensive fields like user interfaces and content services, contrasted by declines in capital-intensive areas such as memory architecture. Valuation heterogeneity is evident, with patent values increasing unevenly across domains, while competitive dynamics show varying competitor numbers over time. To explain these patterns, I build an entry model estimated via moment inequalities that provides bounds on parameters, indicating that a 1\% rise in competitors reduces the firm's innovation profit by 18.44 to 56.85 million USD, an increase of one standard deviation of the portfolio diversification boosts the profit by 0.57 to 36.81 million USD, and the average entry costs are bounded at 75.93 million USD.
% \vspace{0.1in}
% % \section*{Research in Progress}
% \textbf{The Impact of Subsidy on Direction and Scope of Innovation}, with C. Cincotta.
% \bigskip
\section*{Policy Work}
\textbf{The sky in blues: on the recent development of Indonesian airlines industry}. \textit{LPEM FEB UI WP} (2019), with C. Nuryakin, N. W. G. Massiw, A. U. Lumbanraja, S. Hambali, and D. Anindita 
\section*{Miscellaneous}
Language: Indonesian (native), English (fluent), Dutch (basic-intermediate), Arabic (basic-intermediate)

\smallskip
Software: R, Python, Rust
% \bigskip
% Footer
\begin{center}
  \begin{footnotesize}
    Last updated: \today \\
    \href{\footerlink}{\texttt{\footerlink}}
  \end{footnotesize}
\end{center}
\end{document}
